\section{Introduction}
\label{sec:intro}
A macroscopic constitutive law compresses the mechanical behavior of a material into a relation between deformation and stress. For a heterogeneous solid, constructing that relation can be difficult: the arrangement of the constituents, voids, or reinforcing directions determines how the microstructure deforms, and finite deformation can couple loading modes that appear independent at small strains. Computational homogenization retains these mechanisms by computing the macroscopic response from a microscopic boundary-value problem. The imposed macroscopic deformation drives a representative volume element (RVE); microscopic equilibrium then determines the effective stress and the tangent required by the macroscopic solver~\cite{Miehe1999}.

In a two-level finite element calculation, denoted FE$^2$, this microscopic problem is attached to each macroscopic integration point. Its repeated solution during load increments and Newton iterations is the principal computational burden. For a fixed material used in many structural calculations, it is therefore attractive to solve representative microscopic problems offline and reuse the resulting information online. The central choice is \emph{what to approximate}: the microscopic displacement field and its equilibrium equations, or the macroscopic constitutive relation obtained after solving those equations.

These choices lead to projection-based reduced micromechanics and learned constitutive laws, respectively. Both can use artificial neural networks (ANNs), but the role of the network is different. In the first case, it helps represent the microscopic state; in the second, it replaces the effective constitutive response. This paper investigates the accuracy, mechanical structure, and online cost associated with these two approximation routes. Its constitutive contribution addresses a specific difficulty within the second route: how to give a polyconvex network sufficient freedom to represent anisotropy without losing the properties imposed by its construction.

\subsection{Reducing the microscopic problem}
Projection-based model order reduction (PMOR) retains the connection to microscopic equilibrium while reducing the number of unknowns used to describe the displacement field. A conventional projection-based reduced-order model (PROM) approximates that field in an affine space spanned by a reduced-order basis, usually constructed by proper orthogonal decomposition (POD) of previously computed solutions. The coefficients of the basis are the generalized coordinates solved for online. Reduced finite-strain homogenization~\cite{yvonnet2007,Hernandez2014} and potential-based two-scale calculations with GPU acceleration~\cite{Fritzen2016} demonstrate this approach to reducing microscopic cost. Related reduced periodic-substructure methods also operate within domain-decomposition settings~\cite{Hernandez2020}; here, the reduced model is used inside the scale-separated FE$^2$ calculation.

There are two distinct costs to reduce. The first is the dimension of the equilibrium problem. A small affine space may not represent a curved family of microscopic solutions accurately, even when that family depends on only a few loading parameters. Local reduced bases~\cite{Amsallem2012}, autoencoder-based manifolds~\cite{Lee2020}, and quadratic approximation manifolds~\cite{Barnett2022} address this representational difficulty in different ways. The PROM--ANN approach of Barnett et al.~\cite{barnett2023} and the latent-space closure framework of Ares De Parga et al.~\cite{aresdeparga2026nonlinear} instead separate the reduced coordinates into primary and secondary groups. The primary generalized coordinates are retained as independent unknowns; a learned closure predicts the secondary generalized coordinates from them. Thus, a relatively rich displacement representation can be reconstructed while solving for only a few independent coordinates. We use this primary/secondary nomenclature throughout.

The second cost is the evaluation of the reduced equations. A system with few unknowns can remain expensive if its residual requires constitutive evaluations and assembly over the entire microscopic mesh. Sampling-and-weighting hyperreduction addresses this cost by retaining a subset of local contributions and assigning weights that approximate the relevant reduced integrals. Optimized cubature~\cite{An2008}, energy-conserving sampling and weighting (ECSW)~\cite{Farhat2014,farhat2015}, and the empirical cubature method (ECM)~\cite{hernandez2017} provide established constructions, with extensions to local bases and Petrov--Galerkin projection~\cite{Grimberg2021,AresDeParga2023}. The resulting hyperreduced projection-based model (HPROM) therefore combines two approximations: one to the microscopic state and another to the evaluation of its reduced equations.

The integration rule can also exploit how that state representation changes. Continuous empirical cubature adjusts point locations and weights~\cite{Hernandez2024}, whereas subspace-adaptive cubature retains common locations with different weights for different reduced subspaces~\cite{Bravo2024}. Manifold-adaptive weights extend the latter idea to weights that vary continuously with the reduced coordinates~\cite{hernandez2026}. In this study, fixed-weight ECM is used for the affine HPROM and manifold-adaptive cubature for HPROM--ANN. Both are evaluated on actual reduced meshes. They retain an online reduced equilibrium solve; learning a closure does not, by itself, remove that solve. A direct operating mode, D--HPROM--ANN, additionally uses the imposed strain to prescribe strain-informed coordinates and reconstruct the microscopic response without equilibrium iterations. Comparing these modes separates the benefit of a compact representation from the additional approximation introduced by bypassing equilibrium.

\subsection{Learning the macroscopic constitutive response}
Constitutive learning removes the online microscopic problem by approximating its effective response directly. Neural surrogates of homogenized material behavior~\cite{Aldakheel2023} and physics-constrained networks coupled to automated microscopic data generation~\cite{kalina2023} illustrate this alternative. We consider history-independent hyperelasticity, for which the constitutive state can be described by the current deformation. Learning microscopic histories~\cite{Liu2022} or coordinating training across multiple interacting scales in path-dependent materials~\cite{AsadFarhat2026} introduces additional problems beyond this setting.

Even in the history-independent case, fitting stress values is not the only requirement. A macroscopic equilibrium calculation uses both the stress and its derivative with respect to deformation. For a hyperelastic material, these quantities must derive from the same stored-energy function. Directly regressing the stress components does not enforce this relation: a model may fit sampled stresses accurately while predicting nonzero net work around a closed deformation cycle. Learning a scalar energy and obtaining stress by differentiation enforces the potential structure, as emphasized by As'ad et al.~\cite{asad2022}. It also supplies the corresponding major symmetry of the constitutive tangent. These are properties of the represented law, not conclusions drawn from a small fitting error.

An energy potential, however, does not control every aspect of its curvature. This is the motivation for adding convexity constraints, with careful attention to the variables in which convexity is imposed. As'ad et al.~\cite{asad2022} use an input-convex neural network (ICNN) in Green--Lagrange strain. Convexity in this strain measure has a tensile-state stability implication, but it is not the same as polyconvexity in the deformation gradient; the second Piola stress enters the connection between strain-space and deformation-gradient curvature~\cite{GaoNeffRoventaThiel2017}. Polyconvexity instead represents the energy as a convex function of deformation minors~\cite{ball1977}. It provides a tractable sufficient condition for rank-one convexity and, together with the additional growth and admissibility hypotheses, supports existence results in nonlinear elasticity~\cite{Ciarlet1988}. Neither condition should be identified with convergence of a particular finite element solution algorithm.

These distinctions motivate the four constitutive models compared here. Regression predicts stress directly. Free predicts a scalar energy without imposing polyconvexity. The two constrained alternatives use an ICNN or an input-convex Kolmogorov--Arnold network (ICKAN) within a polyconvex energy construction. The comparison first asks whether each model can approximate the same microscopic response, and then examines the additional properties imposed by its architecture. Constitutive tests and structural calculations provide complementary evidence: the former examine the represented law, while the latter measure its consequences in a boundary-value problem.

\subsection{The role of features in polyconvex anisotropic learning}
Anisotropy makes the choice of constitutive inputs consequential. Structural tensors encode preferred material directions, and invariants built from these tensors provide objective measures of deformation. Polyconvex invariant formulations have been developed for anisotropic elasticity~\cite{SchroderNeff2003}, while models of distributed fiber orientations illustrate the importance of representing directional structure beyond a single aligned family~\cite{Gasser2006}. In a neural constitutive model, these scalar deformation measures become the network's \emph{features}.

One sufficient neural construction combines polyconvex features with an energy core that is convex and nondecreasing in each feature. Input-convex neural networks provide the convex core~\cite{amos2017}; anisotropic invariant formulations and compatible reference-state corrections are developed by Klein et al.~\cite{klein2022} and Linden et al.~\cite{linden2023}. Other structured energy parameterizations include constitutive artificial neural networks~\cite{Linka2023} and polyconvex neural ordinary differential equations~\cite{Tac2022}. For spline-based Kolmogorov--Arnold networks (KANs), the distinction between constitutive KANs and input-convex KANs is important: the CKAN framework combines energy modeling, partial monotonicity, and symbolic interpretation~\cite{Abdolazizi2025}, whereas the monotonic ICKAN construction adds the convexity constraints needed for its polyconvex invariant composition~\cite{thakolkaran2025}. Here, ICNN and ICKAN are used as alternative cores within a common feature design.

The features are not a neutral preprocessing step. They restrict the functions that the constrained core can represent. For example, if one deformation produces a feature vector that is larger componentwise than another, a nondecreasing core cannot assign it a lower output. Increasing the number of neurons does not remove this ordering restriction. Reference corrections must also be accounted for when considering the complete energy, but the underlying representational constraint remains. Klein et al.~\cite{klein2022} provide a pertinent numerical example: their invariant-based model struggles with shear in a cubic lattice, whereas a deformation-gradient-based polyconvex model gives a better fit. The invariant approach performs well on their transversely isotropic example. Poor approximation by one admissible representation therefore does not, on its own, establish an incompatibility between the material response and polyconvexity.

The resulting design problem is to enlarge the feature representation while retaining the conditions used to establish admissibility. Learning anisotropy itself is already addressed through generalized structure tensors~\cite{Kalina2024}, and objective triclinic polyconvex formulations can be constructed without prescribing a nontrivial material symmetry group~\cite{klein2026}. Our focus is more specific: paired directional features with trainable orientations, stretch powers, and determinant exponents constrained by explicit joint-convexity bounds. The pairing is chosen so that each feature derivative at the undeformed state is proportional to the identity tensor. The resulting reference stress can therefore be cancelled by a single scalar correction affine in the determinant, leaving the positive volumetric stiffness terms independently adjustable. Away from the reference state, the paired terms can still produce an anisotropic response. The approach connects representational flexibility to a reference normalization that preserves the polyconvex construction.

\subsection{Contribution and assessment}
The first contribution is this paired-feature constitutive construction and its realization with a smooth ICNN and an integrated-hat spline ICKAN. The feature family and reference normalization are shared, so the two cores can be assessed within the same constitutive framework. The construction establishes objectivity, a stress-free reference, two-dimensional polyconvexity, and specified growth behavior. Nonnegative energy is addressed separately by a sufficient lower bound checked for the selected trained parameters. The paired features and analytically matched spline tails are the specific ingredients developed here; the underlying ICNN, ICKAN, and affine-determinant normalization principles build on the cited work.

The second contribution is a common assessment of these learned energies against reduced micromechanics and nested full-order FE$^2$ (FOM--FE$^2$). The affine HPROM, the iterative HPROM--ANN, and the direct decoder mode provide different levels of retained microscopic information and online equilibrium enforcement. Their comparison establishes the accuracy--cost alternatives available for the same constitutive task, including the effect of retaining or bypassing reduced equilibrium. POD, latent-space closure, ECM, and manifold-adaptive cubature are established ingredients of this comparison. Their constitutive consequences depend on the actual projection, integration, and output evaluation; they do not automatically inherit the macroscopic polyconvexity certificate of a separately constructed energy network. Figure~\ref{fig:hierarchy} summarizes the two approximation routes.

The numerical examples separate constitutive assessment from structural deployment. The periodic-cell study tests whether the constrained feature family can fit an anisotropic homogenized response, then uses closed-cycle work and rank-one curvature to examine properties not established by pointwise stress accuracy. This attention to derivatives also follows the constitutive benchmarking of Tac et al.~\cite{Tac2024}. The tensile-coupon study uses the same material to measure how approximation errors propagate into structural fields and complete online wall time, with nested FOM--FE$^2$ as the reference. Field comparisons supplement aggregate errors, as also motivated by reduced-model assessments such as Ref.~\cite{Chmiel2024}. The two examples thus answer different questions without representing the coupon as a second microstructure. The present evidence concerns one periodic material, plane-strain kinematics, and a specified deformation domain.

Section~\ref{sec:fom} defines the homogenization problem and work-conjugate variables. Section~\ref{sec:reduced} presents the reduced microscopic models, and Section~\ref{sec:learned} develops the learned constitutive laws. Sections~\ref{sec:cell} and \ref{sec:coupon} report the constitutive and structural assessments, respectively.

\begin{figure}[htbp]
 \centering\includegraphics[width=.98\linewidth]{model_hierarchy.pdf}
 \caption{Two approximation routes for a macroscopic constitutive evaluation. Reduced micromechanics reconstructs a microscopic state and evaluates the microscopic law; a learned constitutive energy supplies the effective stress directly by differentiation. The iterative HPROMs retain a reduced equilibrium solve, while the direct decoder mode bypasses it.}
 \label{fig:hierarchy}
\end{figure}
